\begin{theorem}[椭圆平均的径向留数测度：唯一性与可审计反演]\label{prop:xi-singular-ring-ellipse-wallcrossing-atomic-measure}
沿用定理 \ref{thm:xi-singular-ring-ellipse-wallcrossing-rational} 的记号，并令 \(r=\sqrt{L}>1\)。
定义有限原子复测度
\[
\mu_f:=\sum_{\substack{a\neq 0\\ a\ \text{为 }h\text{ 的极点}\\ |a|>1}}\Res_{z=a}h(z)\,\delta_{|a|},
\]
其中 \(\delta_s\) 为点 \(s\in(1,\infty)\) 处的 Dirac 质量。
则对任意 \(1<r_1<r_2\)，成立壁穿越差分恒等式
\[
\mathcal{L}_{r_2^2}(f)-\mathcal{L}_{r_1^2}(f)=\mu_f\bigl((r_1,r_2)\bigr).
\]
并且 \(\mu_f\) 由上述差分关系唯一确定；更进一步，
\[
\mathcal{L}_{r^2}(f)=\mathcal{L}_{1^+}(f)+\mu_f\bigl((1,r)\bigr),
\]
从而函数 \(r\mapsto \mathcal{L}_{r^2}(f)\) 与径向测度 \(\mu_f\)（外加积分常数 \(\mathcal{L}_{1^+}(f)\)）之间一一对应。
\end{theorem}

\begin{proof}
由定理 \ref{thm:xi-singular-ring-ellipse-wallcrossing-rational}，当 \(r_1,r_2\) 均不等于 \(h\) 的任一非零极点模长时，
\[
\mathcal{L}_{r_i^2}(f)
=\Res_{z=0}h(z)\;+\;\sum_{\substack{a\neq 0\\ a\ \text{为 }h\text{ 的极点}\\ 0<|a|<r_i}}\Res_{z=a}h(z)
\qquad(i=1,2).
\]
相减得到
\[
\mathcal{L}_{r_2^2}(f)-\mathcal{L}_{r_1^2}(f)
=\sum_{\substack{a\neq 0\\ a\ \text{为 }h\text{ 的极点}\\ r_1<|a|<r_2}}\Res_{z=a}h(z)
=\mu_f\bigl((r_1,r_2)\bigr).
\]
对一般的 \(1<r_1<r_2\)，由两侧皆为分段常值函数且仅在离散集合处发生跳跃，可取不穿越极点模长的逼近序列并由一致性推出同一差分恒等式。
\par
若 \(\mu_f'\) 亦满足同一差分关系，则对任意区间 \((r_1,r_2)\) 有 \(\mu_f((r_1,r_2))=\mu_f'((r_1,r_2))\)，由原子测度在该区间族上的一致性推出唯一性。
最后令 \(r_1\downarrow 1\) 即得
\(
\mathcal{L}_{r^2}(f)=\mathcal{L}_{1^+}(f)+\mu_f((1,r))
\)，从而建立可审计的反演关系。证毕。
\end{proof}

\endinput

